\notatka[04.11.2026]{Przykładowa notatka}

Zwykły tekst z \pojecie{ważnym pojęciem} i wzorem $E = mc^2$, a temperatura to $95\degree$C.

\begin{fakty}[Fakty o projekcie]
Koszt: 3 miliardy dolarów\\
Czas trwania: 15 lat
\end{fakty}

\begin{definicja}[PCR]
Reakcja łańcuchowa polimerazy namnaża wybrany fragment DNA.
\end{definicja}

\begin{twierdzenie}[Przykład]
Wzory i twierdzenia dzielą numerację w obrębie sekcji.
\end{twierdzenie}

\begin{uwaga}
Coś, o czym warto pamiętać.
\end{uwaga}

\begin{wazne}[Na kolokwium]
ddNTP nie ma grupy \mbox{-OH} na końcu 3$'$.
\end{wazne}

\subsection{Podsumowanie}
\begin{podsumowanie}
Jedna notatka to jeden plik w \texttt{content/}.
\end{podsumowanie}
